%!TEX root = ../thesis.tex
% Chapter --- Background

\chapter{Background} \label{sec:background}

This chapter builds the ideas the project rests on: why texture features suit mammograms, how the wavelet scattering transform turns texture into a stable representation, and --- the pivot of the whole project --- why the way that representation is \emph{summarised} decides whether it works on the tissue task or the mass task. Convolutional neural networks and their fusion with scattering features, the other half of the replicated pipeline, come last. Implementation choices and parameters are deferred to Chapter~\ref{sec:methodology}.

\section{Texture and the two classification tasks}

A radiologist reads each mammogram and records findings on the American College of Radiology's BI-RADS scale, which assigns both an overall suspicion category and a breast-density category \cite{sickles2013acr}; these two axes correspond directly to the two classification problems this project addresses, following Razali et al.\ \cite{razali2023cnn}.

\textbf{Tissue-density classification} decides whether a region of breast tissue is fatty or dense (fibroglandular). Density matters clinically because dense tissue both raises cancer risk and hides tumours behind overlapping tissue \cite{boyd2007mammographic}, and it is essentially a property of \emph{relatively uniform texture} --- the overall pattern of a region rather than any single structure within it. \textbf{Mass classification} decides whether a detected lesion is benign or malignant, and is a different kind of problem: the answer lies in a \emph{localised, structured object}, whose diagnostic cues are the shape of its margin and the arrangement of its interior, not a uniform texture. This texture-versus-structure distinction is the thread running through the whole report and, as this chapter argues, it dictates how the scattering representation should be used.

Both tasks turn on \emph{texture}: the spatial pattern of intensity variation rather than overall brightness. A classical way to quantify texture is the grey-level co-occurrence matrix \cite{haralick1973textural}, which tallies how often given pairs of grey levels occur at a fixed offset --- but it captures a single scale, whereas breast structure appears at many sizes at once, from fine trabecular strands to broad glandular regions. A representation that measures structure across \emph{multiple scales} is therefore better matched to the problem, which is what motivates wavelet-based features. Such handcrafted features are especially attractive here because they are computed by a fixed rule: they need no training data, are deterministic, easy to interpret, and cannot overfit --- decisive when annotated mammograms are scarce.

\section{Wavelets and the scattering transform} \label{sec:wst-bg}

A wavelet is a small, localised oscillating waveform \cite{morlet1982wave}. Dilating and rotating a single ``mother'' wavelet produces a bank of filters, each tuned to a particular scale and orientation; representing an image through such a multi-scale filter bank is the basis of wavelet image analysis \cite{mallat1989theory}. Convolving an image with this bank measures how much oscillatory, edge-like content it contains at each scale, orientation and location: coarse wavelets respond to large, smooth structures and fine wavelets to small, sharp ones. Because each wavelet is localised in space --- unlike the global sine waves of the Fourier transform --- the response changes only slightly when the image is slightly warped. Its shortcoming is that the raw response is \emph{not} translation-invariant: shift the image and every coefficient shifts with it, so two patches of the same tissue at different positions look different to a classifier.

The wavelet scattering transform (WST), introduced by Mallat and by Bruna and Mallat \cite{mallat2012invariant, bruna2013invariant}, removes this nuisance sensitivity while keeping the useful texture information. It repeats three simple steps. First it \emph{convolves} the image with an oriented wavelet. Second it applies the \emph{modulus} (absolute value), which strips out the oscillating sign of the response and leaves a smooth, positive ``envelope''; this non-linear step is essential, because a wavelet averages to zero, so without it the next step would cancel the signal out entirely. Third it \emph{averages} the envelope with a low-pass filter at scale $2^J$, and this averaging is exactly what makes the descriptor invariant to shifts smaller than $2^J$. Averaging inevitably discards fine detail, but that detail is not lost: feeding the envelope back through another wavelet-and-modulus step recovers it, and repeating the process builds up successive \emph{scattering orders} $S_0, S_1, S_2, \dots$ --- the zeroth order being simply the blurred image, and each higher order recovering structure that the previous averaging discarded. The full definitions are given in Section~\ref{sec:ws}.

Three properties make this representation well suited to small mammography datasets. It is \emph{locally translation-invariant}, so the same tissue at different positions maps to nearly the same features. It is \emph{stable to small deformations}: a small warp of the image changes the coefficients only a little, and in a controlled, bounded way (a Lipschitz-continuity guarantee \cite{mallat2012invariant}), so the natural variation between patches of one class does not confuse the descriptor. And almost all of its energy is captured within the first two orders, so second-order scattering is the standard and sufficient choice \cite{bruna2013invariant}. Crucially, the wavelet filters are fixed in advance: the WST has nothing to learn and therefore cannot overfit, which is why it can be applied directly to a few hundred patches.

\section{How the representation is summarised: the central choice}

Each of the transform's outputs is a small coefficient \emph{map} --- one image per ``path'', a path being one particular sequence of scales and orientations through the cascade. The design choice at the heart of this project is how to reduce these maps to a single feature vector, and there are three options. \emph{Spatial averaging} replaces each map by its mean, giving one number per path; this is the most compact and most invariant summary, and it treats the whole patch as one homogeneous texture. \emph{Preserving the maps} instead keeps the spatial grid, retaining information about \emph{where} structure sits, at the cost of some translation-invariance. A middle option, the \emph{scattering covariance} of Cheng and M\'enard \cite{cheng2021scattering}, keeps the averaged coefficients but adds correlations \emph{between} different paths, enriching the descriptor while it remains a compact set of scalars. The reference paper, like the original scattering work, uses plain spatial averaging \cite{bruna2013invariant, razali2023cnn}.

Whether averaging is the right choice depends entirely on the task. For tissue density it is ideal: the discriminative signal is the \emph{statistics} of a roughly uniform texture, so collapsing each map to its mean throws away only the irrelevant position information. For a mass it is damaging. Malignancy is written into the spatial arrangement of distinct sub-regions --- a dense \emph{core}, the \emph{margin} (spiculated or irregular margins point to malignancy, smooth well-defined ones to benignity), and the immediately \emph{surrounding} tissue, which may show architectural distortion --- and averaging over the whole patch blends these regions into a single value, erasing the very core-versus-margin contrast that separates benign from malignant. This is exactly the pattern the results confirm (Chapter~\ref{sec:results}): averaged scattering is a strong \emph{tissue} descriptor but a weak \emph{mass} one. The project's two extensions follow directly from this single observation --- a richer \emph{scalar} descriptor (scattering covariance) for tissue, and \emph{preserving the spatial map} for masses (Chapter~\ref{sec:extensions}).

\section{Learned features and their fusion with scattering} \label{sec:cnns-bg}

A convolutional neural network (CNN) is the learned counterpart to these handcrafted features: like the WST it slides banks of filters across the image, but its filters are \emph{learned} from labelled data by gradient descent rather than fixed in advance \cite{lecun1998gradient, he2016deep}. This lets a CNN discover whatever features best separate the classes, but at a price --- it needs a large amount of labelled data to do so and tends to overfit when data is scarce, which is precisely the situation for mammography. The two feature types are therefore complementary: the WST is fixed and robust but cannot adapt, while the CNN is adaptable but data-hungry. This is the motivation for \emph{fusion} --- combining them into a single decision. The reference paper concatenates the CNN and scattering feature vectors and classifies the combined vector \cite{razali2023cnn}, so that the stable, training-free scattering features can support the CNN exactly when it is data-starved and most at risk of overfitting.


% ============================================================
% 3. METHODOLOGY
% ============================================================
